\section{The Supervised Fit at Increasing Budgets}\label{sec:attack:budget}

The supervised fit meets the criterion of H2.
At \vf{1{,}000} paired sequences on Corpus A, word-F1 is \vf{0.876}, against a hub reference of \vf{0.946} and a mismatched-prompt control of \vf{0.11}; half of the reference is \vf{0.473}. % src: results/v7_2/main/mixed-v7_9-s4.json seed_summary.budgets.1000.macro_f1.b1000/a1/r1 (0.8759), seeds.*.cells.b1000/a1/r1.wrong_memory (0.1087 to 0.1154); factorial.json controls.true_olmo_code (0.9459)
Fitted on \vf{1{,}000} prompts, an affine map thus brings the hidden states of Pythia within \vf{0.07} word-F1 of the hub's own OLMo outputs.
The scores are means over three seeds on the \vf{500} selection prompts.

Word-F1 rises with the budget at every step (Table~\ref{tab:attack:budget}, Figure~\ref{fig:attack:budget}).
The supervised fit scores \vf{0.348} at \vf{10} paired sequences and \vf{0.758} at \vf{100}. % src: mixed-v7_9-s4.json seed_summary.budgets.{10,100}.macro_f1.b*/a1/r1
The knee is the smallest budget whose gain above the control reaches \vf{95} percent of the gain at the largest budget.
The knee is \vf{300}: the fit at \vf{300} recovers \vf{96} percent of the gain at \vf{1{,}000}, and the fit at \vf{100} recovers \vf{85} percent. % src: computed from Table values and the controls of the same cells (0.8452 - 0.1103) / (0.8759 - 0.1125) = 0.963; (0.7582 - 0.1102) / 0.7634 = 0.849
Budgets of \vf{10} to \vf{100} paired sequences are therefore where the supervised fit leaves most to recover, and the following sections concentrate on them.

\begin{table}[H]
\caption{The supervised fit rises from \vf{0.35} at \vf{10} paired sequences to \vf{0.88} at \vf{1{,}000}, against a hub reference of \vf{0.946}.
Word-F1 averaged over three decoders and three seeds, Pythia on Corpus A, \vf{500} selection prompts; the range over seeds is at most \vf{0.051} in every cell.
The mismatched-prompt control of these cells lies between \vf{0.10} and \vf{0.15}.}
\label{tab:attack:budget}
\centering
\begin{tabular}{lccccc}
\toprule
Budget (paired sequences) & 10 & 30 & 100 & 300 & 1{,}000 \\
\midrule
Supervised fit & \vf{0.348} & \vf{0.536} & \vf{0.758} & \vf{0.845} & \vf{0.876} \\ % src: results/v7_2/main/mixed-v7_9-s4.json seed_summary.budgets.*.macro_f1.b*/a1/r1.mean (0.3482, 0.5355, 0.7582, 0.8452, 0.8759)
Supervised fit without the ridge penalty & \vf{0.303} & \vf{0.441} & \vf{0.656} & \vf{0.817} & \vf{0.884} \\ % src: same file, b*/a1/r0.mean (0.3028, 0.4412, 0.6558, 0.8174, 0.8837); seed range at most 0.051 (b10/a1/r0)
\bottomrule
\end{tabular}
\end{table}

\begin{figure}[ht]
\centering
\includegraphics[width=0.75\textwidth]{ch4/budget_curve}
\caption{The hybrid fit exceeds the supervised fit by \vf{0.31} word-F1 at \vf{10} paired sequences and by \vf{0.01} at \vf{300}.
Word-F1 averaged over three decoders on the \vf{500} selection prompts of Corpus A; lines are means over three seeds and bands the range over seeds.
The dotted line is the hub reference, and the gray band is the mismatched-prompt control of the plotted fits.}
\label{fig:attack:budget}
\end{figure}

The ridge penalty raises the supervised fit at \vf{10} to \vf{300} paired sequences and lowers it by \vf{0.008} at \vf{1{,}000} (Table~\ref{tab:attack:budget}).
The gain at small budgets is consistent with the size of the map: about one million weights, against \vf{92} rows at a budget of \vf{10}.

The three decoders do not read the mapped states equally.
At \vf{1{,}000} paired sequences the supervised fit scores \vf{0.964} through SmolLM, \vf{0.886} through OLMo, and \vf{0.778} through Qwen, against hub references of \vf{0.986}, \vf{0.971}, and \vf{0.881}. % src: mixed-v7_9-s4.json seeds.*.cells.b1000/a1/r1.final.per_decoder, mean of three seeds (0.9643, 0.8855, 0.7780); factorial.json controls.true_olmo_code.per_decoder (0.9855, 0.9714, 0.8806)
The three-decoder mean therefore hides a distance to the reference of \vf{0.02} for SmolLM and \vf{0.10} for Qwen.

The result does not show how the fit behaves on held-out prompts or on a second corpus: all scores come from the selection prompts of Corpus A.
